\documentclass[11pt]{article}

\usepackage[margin=1in]{geometry}
\usepackage[T1]{fontenc}
\usepackage[utf8]{inputenc}
\usepackage{lmodern}
\usepackage{microtype}
\usepackage{booktabs}
\usepackage{array}
\usepackage{xcolor}
\usepackage{hyperref}
\usepackage{listings}
\usepackage{enumitem}

\hypersetup{
  colorlinks=true,
  linkcolor=blue!55!black,
  citecolor=blue!55!black,
  urlcolor=blue!55!black
}

\definecolor{codebg}{RGB}{248,248,250}
\definecolor{codeframe}{RGB}{210,210,216}

\lstset{
  basicstyle=\ttfamily\small,
  backgroundcolor=\color{codebg},
  frame=single,
  framerule=0.4pt,
  rulecolor=\color{codeframe},
  breaklines=true,
  keepspaces=true,
  showstringspaces=false,
  xleftmargin=6pt,
  xrightmargin=6pt
}

\title{FSL: A Governed Symbolic Language and Formal Observer Theorem Package}
\author{Simon DeFrisco\\Independent Researcher\\\texttt{formatlifeofficial@gmail.com}}
\date{Public package v1.1.3}

\begin{document}
\maketitle

\begin{abstract}
FSL is a governed symbolic language for publishing, evolving, and verifying theorem-facing semantics in agent systems. This paper presents the public FSL observer theorem package: a release containing the FSL language snapshot, observer theorem registry, Lean source snapshots, theorem lifecycle export, assumptions audit, checksum manifest, and formal proof-status whitepaper. The package contains 32 public theorem records. The current formal coverage audit classifies 31 records as machine-checked theorem records and 1 record as machine-checked under explicit cryptographic axioms. No records are partial, definition-only, or planned. The contribution is not an assumption-free proof of all system boundaries; it is an evidence-disciplined publication package that separates checked theorem claims, axiom-dependent claims, lifecycle attestations, generated language exports, runtime parity evidence, and excluded claims.
\end{abstract}

\section{Introduction}

Formal systems for AI governance and agent coordination often mix several kinds of claims: symbolic notation, executable rules, theorem statements, runtime checks, cryptographic evidence, and publication metadata. FSL was designed to keep those layers separate while still allowing them to be published together.

The FSL observer theorem package is a public release boundary for that design. It packages a governed symbolic language with a bounded-observer theorem stack: agents have cryptographic identity, bounded spatial position, monotone temporal state, finite observation horizons, and proof-bearing commitments. The public release is intended for readers who need an inspectable artifact without access to the private SiMON source repository.

This PDF is the reader-facing paper layer. The authoritative release artifacts are the repository files:

\begin{itemize}[leftmargin=*]
  \item \texttt{formal\_whitepaper.md}, the theorem-by-theorem formal proof-status paper;
  \item \texttt{FORMAL\_PROOF\_BUNDLE.md}, the proof bundle guide;
  \item \texttt{THEOREM\_REGISTRY.md} and \texttt{theorem\_registry.json}, the public theorem registry;
  \item \texttt{theorem\_lifecycle.json}, the public lifecycle export;
  \item \texttt{lean\_coverage\_report.md} and \texttt{lean\_coverage.json}, the Lean coverage audit;
  \item \texttt{ASSUMPTIONS\_APPENDIX.md} and \texttt{lean\_assumptions.json}, the assumptions and proof-hole audit;
  \item \texttt{fsl/SYSTEM.yaml}, the canonical FSL source snapshot;
  \item \texttt{CHECKSUMS.sha256}, the public package integrity file.
\end{itemize}

\section{Package Claim}

The public package makes a deliberately narrow claim:

\begin{quote}
The FSL observer theorem package contains 32 public theorem records with explicit proof-status evidence: 31 machine-checked theorem records and 1 machine-checked-under-explicit-axioms record.
\end{quote}

The package does not claim that every assumption has been eliminated, that SHA-256 collision resistance is proved inside Lean, that Rust runtime surfaces are constitutional authority, or that lifecycle attestation is itself a Lean proof.

\section{FSL as Publication Language}

FSL is a governed symbolic language. Its public release includes canonical identifiers, symbol metadata, type and family classifications, theorem aliases, compatibility metadata, and grammar/rendering exports. The language layer matters because theorem packages need stable public names. Human labels may change; canonical identifiers and theorem references should not.

In this release, \texttt{fsl/SYSTEM.yaml} is the canonical FSL source snapshot. The JSON and Markdown exports are derived publication artifacts. They are included for readability and external tooling, but they do not replace the canonical source snapshot.

\section{Observer Theorem Model}

The observer model treats an agent as a bounded state:

\begin{lstlisting}
observer_state = (DID, cell, depth, tick)
\end{lstlisting}

\noindent where \texttt{DID} is a cryptographic identity, \texttt{cell} and \texttt{depth} locate the observer in a Hierarchical Triangular Mesh, and \texttt{tick} is a monotone temporal coordinate. Lawful motion preserves identity, advances time, and restricts spatial motion to admissible transitions.

The theorem stack covers HTM geometry, observer decomposition, motion, horizon bounds, cost bounds, ambiguity, commitment, impossibility results, and FSL movement records.

\section{Formal Coverage}

Table~\ref{tab:coverage} summarizes the current public proof-status boundary.

\begin{table}[h]
\centering
\begin{tabular}{lr}
\toprule
Classification & Count \\
\midrule
Machine-checked theorem records & 31 \\
Machine-checked under explicit axioms & 1 \\
Partial machine-checked statements & 0 \\
Definition-only records & 0 \\
Planned / not yet encoded records & 0 \\
\midrule
Total public theorem records & 32 \\
\bottomrule
\end{tabular}
\caption{Formal coverage summary for public package v1.1.3.}
\label{tab:coverage}
\end{table}

The axiom-dependent theorem is \texttt{gbo\_vi\_non\_equivocating}. It is checked in Lean under an explicit cryptographic commitment boundary. The boundary is documented in \texttt{CRYPTO\_AXIOM\_BOUNDARY.md}.

\section{Assumptions and Trusted Boundaries}

The public assumptions audit scanned 10 Lean files. It found 132 declared axioms and 0 code-level \texttt{sorry} or \texttt{admit} proof holes. The largest assumption surface is the governance model and bridge file, which declares propositional governance atoms and bridge relations. The cryptographic commitment theorem depends on four model primitives:

\begin{itemize}[leftmargin=*]
  \item \texttt{Commitment}
  \item \texttt{commit}
  \item \texttt{ComputationallyBinding}
  \item \texttt{sha256\_binding\_assumption}
\end{itemize}

These assumptions are not hidden. They are part of the claim boundary. The theorem \texttt{gbo\_vi\_non\_equivocating} is therefore a checked theorem under an abstract commitment and binding model, not an assumption-free proof of SHA-256 or any concrete cryptographic implementation.

\section{Lifecycle Evidence}

The theorem lifecycle export reports 32 active theorem records and a valid lifecycle chain. Lifecycle status is publication evidence: it records that a theorem record has been attested into the governed theorem lifecycle. It is not a substitute for Lean proof status.

This distinction is central to the package. Lean proof status, lifecycle status, FSL registry status, checksum status, and runtime parity status are separate evidence layers.

\section{Artifact Availability}

The public repository is intended to be the citable artifact for this work:

\begin{quote}
\url{https://github.com/ProphetGang/formal_symbol_language}
\end{quote}

The repository should be cited by version tag after release:

\begin{quote}
\texttt{fsl-observer-theorem-v1.1.3}
\end{quote}

A Zenodo DOI may be added after the GitHub release is archived. Until then, the DOI field should be treated as pending:

\begin{quote}
DOI: pending Zenodo archive
\end{quote}

\section{Verification}

From the public package directory, readers can verify package integrity with:

\begin{lstlisting}
shasum -a 256 -c CHECKSUMS.sha256
\end{lstlisting}

The public package intentionally excludes private keys, private governance state, local caches, model files, Lean build products, embedded repositories, and private runtime material.

\section{Conclusion}

The FSL observer theorem package is a formal proof-status release, not merely a narrative paper. It publishes the language, theorem registry, proof-status audit, lifecycle export, assumptions appendix, Lean source snapshots, and checksums required to inspect the package boundary. The result is an evidence-disciplined public artifact: every theorem claim is tied to a status, every known axiom boundary is named, and every excluded claim is made explicit.

\bibliographystyle{plain}
\bibliography{references}

\end{document}
